\chapter{Installation}

A valid installation looks like this: The \linnet{} application is a
simple executable file plus a folder with resources for the Octave code
generator; this folder is called \file{private} and is located next to the
\linnet{} executable file, i.e., both are siblings in the same parent
directory. The environment variable \code{LINNET\_HOME} points to this
parent directory.

Accordingly, this is what you have to do to install \linnet{}:

Create an empty folder, the installation directory, e.g.,
\file{c:\bs{}ProgramFiles\bs{}linNet} (Windows) or \file{/usr/bin/linNet}
(Linux, Mac OS). Copy the \linnet{} executable file and the folder
\file{private} into this directory. Create a persistent environment
variable \code{LINNET\_HOME} and assign the path to the installation
directory.

Linux and Mac OS users need to open a shell and run command \code{chmod +x
/usr/bin/linNet/linNet} to make the binary file executable.

The usual way to use \linnet{} is to run it from a shell window. You will
probably organize your input circuits in a folder or a set of such and
\code{cd} to this folder. To run \linnet{} in this scenario in a
convenient way, it is useful to place a reference to the installation
directory in the system environment variable \code{PATH}. In our example
you'd append \code{;\bs{}ProgramFiles\bs{}linNet} (Windows) or
\code{:/usr/bin/linNet} (Linux, Mac OS) to the current value of this
variable.

The definition of a shell command-alias pointing to the \linnet{}
executable file can be a reasonable alternative to extending the
environment variable \code{PATH}.

If environment variable \code{LINNET\_HOME} is not set then \linnet{}
makes an attempt to find its resources via the path to itself. However, it
depends on the system, the shell in use and the way the command is typed,
whether this information is completely passed to the application or not.
If it can't find out where to find itself it'll refuse to generate Octave
code as it lacks the folder \file{private} with code
templates.\footnote{There are system specific functions which let an
application find out, where its executable file is located, e.g.,
\code{GetModuleFileName} for Windows, but the use of such functions would
make the source code of \linnet{} platform dependent. This price is too
high for a slightly simpler installation.} Not setting the environment
variable can be a safe option if you decide to run \linnet{} solely from a
shell window with an alias, which expands to the full, absolute path of
the executable.

\section{Windows}

\subsection{Pre-built executables}

% TODO Add link to GitHub Releases page

The download of \linnet{} contains pre-build, ready-to-use executable
files for Windows (DEBUG and PRODUCTION compilation). The executable files
are already bundled with the resources folder \file{private} so that the
installation reduces to copying the wanted build folder -- \file{DEBUG}
or \file{PRODUCTION} -- to any location of choice, preferably
\file{c:\bs{}\-Program\-Files}, renaming the folder itself to
\file{\linnet{}} and setting the environment variable \code{LINNET\_HOME};
it would be set to \file{c:\bs{}\-Program\-Files\bs{}\-\linnet{}} if the
suggested installation folder is used. The chosen installation path
should not contain blanks.

\subsection{Required DLLs}

The build of \linnet{} has the executable file \file{linNet.exe} as only
product. This executable file is however not self-contained; the
application is linked against the compiler's C libraries, which in case of
the GNU compiler collection build on some Windows DLLs. These DLLs need to
be available in order to execute \linnet{}. Normally, this should always
be the case. In the rare case that such a library should be missing on
your system then the only way out will be to install the GCC port of MinGW
and to rebuild \linnet{}. This would bring all required DLLs into your
computer -- if you can run the simplest \code{helloWorld} application
compiled with GCC then you should also be able to run \linnet{}.

The 64 Bit build of \linnet{} with GCC \ident{mingw-w64-gcc-13.2.0-x64} on
a Windows 11 system loads these DLLs:

\begin{verbatim}
C:\WINDOWS\SYSTEM32\ntdll.dll
C:\WINDOWS\System32\KERNEL32.DLL
C:\WINDOWS\System32\KERNELBASE.dll
C:\WINDOWS\System32\ucrtbase.dll
\end{verbatim}

Other target platforms or other compiler releases will probably require
other DLLs. For example, the 64 Bit build with GCC
\ident{mingw-w64-gcc-mcf\_20260501\_16.1.0\_x64-msvcrt}, also on Windows
11, requires an additional DLL, which is located in the compiler's
installation directory:

\begin{verbatim}
libmcfgthread-2.dll
\end{verbatim}

To run the \linnet{} executable in this scenario, you would either have to
copy that particular DLL into the folder, where \file{linNet.exe} is
located or you would need to add folder \file{mingw64\bs{}bin} in the
installation of GCC to the system environment variable \code{PATH}, so
that Windows can find this DLL.

\subsection{Start from Windows Explorer (file name association)}

If you want to start the computation of transfer functions directly from
the Windows Explorer then you can add a file class definition to the
Windows Registry. It'll extend the context menu entries of the Explorer
for circuit netlist files. With right-click on the file you can load the
netlist file into a text editor and by double-click you can run the
computation with \linnet{}. Please note, if starting \linnet{} this way
Windows will swallow all the program output and you will have to open the
log file after the computation to find out if the computation was
successful.

Here's a template for a Windows registry script to make an appropriate
file class definition:

% TODO Double-check appropriateness of \clearpage command after any change of chapter Installation
\clearpage
{\footnotesize
\begin{verbatim}
REGEDIT4

[HKEY_CURRENT_USER\Software\Classes\.cnl]
@="FileClass_cnl"

[HKEY_CURRENT_USER\Software\Classes\.ckt]
@="FileClass_cnl"

[HKEY_CURRENT_USER\Software\Classes\FileClass_cnl]
@="linNet Circuit Netlist"

[HKEY_CURRENT_USER\Software\Classes\FileClass_cnl\shell]
@="Edit"

[HKEY_CURRENT_USER\Software\Classes\FileClass_cnl\DefaultIcon]
@="c:\\ProgramFiles\\linNet\\linNet.exe,1"

[HKEY_CURRENT_USER\Software\Classes\FileClass_cnl\shell\Edit]
@="&Edit"
[HKEY_CURRENT_USER\Software\Classes\FileClass_cnl\shell\Edit\command]
@="c:\\ProgramFiles\\emacs-30.2\\bin\\emacsclient.exe -n \"%1\""

[HKEY_CURRENT_USER\Software\Classes\FileClass_cnl\shell\Compute]
@="Run linNet"
[HKEY_CURRENT_USER\Software\Classes\FileClass_cnl\shell\Compute\command]
@="c:\\ProgramFiles\\linNet\\linNet.exe -l -c -s -o -- \"%1\""
\end{verbatim}
} % End of decreased script size for the listing

Copy this code into a text file. Save it using a file name with extension
\file{.reg}, e.g., \file{fileClassCnl.reg}.

Please note, that you will have to modify the proposed registry code as it
uses absolute paths to the installed applications; replace these paths
with paths from your computer. Furthermore, the text editor for editing
the circuit netlists is just a suggestion; you will probably replace the
reference to Emacs with a reference to your favorite text editor.

After customization, double-click this file. If your system has no file
name association for \file{*.reg} files then you need to run the registry
script from a shell window. Submit a command like \code{regedit.exe
fileClassCnl.reg}. Administrator rights are not required to complete.

\section{Linux and Mac OS}

Unfortunately, Linux and Macintosh users will have to build \linnet{} from
the sources prior to use. At the moment no pre-built executables are
available for these systems. Please refer to section
\ref{secBuildingLinNet} on page \pageref{secBuildingLinNet}.
